% chapters/07/latihan.tex
\bagian{Latihan dan Refleksi}

\begin{enumerate}
  \item Jelaskan mengapa benda permainan sebaiknya disusun sebagai kelas
        turunan \code{pygame.sprite.Sprite}, bukan disimpan pada beberapa
        variabel terpisah. Sebutkan tiga alasannya.

  \item Sebutkan dua nama anggota yang wajib dimiliki sebuah \istilah{sprite}
        pada Pygame, dan jelaskan apa yang terjadi apabila namanya berbeda.

  \item Sebutkan tiga perintah pemeriksaan tabrakan yang tersedia pada Pygame,
        lalu jelaskan keperluan masing-masing.

  \item Jelaskan perbedaan antara mendeteksi tabrakan dan menyelesaikan
        tabrakan. Mengapa yang pertama dapat diserahkan kepada pustaka,
        sedangkan yang kedua tidak?

  \item Jelaskan mengapa penyelesaian tabrakan dilakukan sumbu demi sumbu,
        bukan sekaligus. Sebutkan gejala yang muncul apabila kedua sumbu
        disatukan.

  \item Sebuah benda bergerak 800 piksel per detik pada laju 40 bingkai per
        detik. Hitunglah perpindahannya pada satu bingkai, lalu tentukan tebal
        bongkahan paling kecil yang masih dapat menghalanginya.

  \item Jelaskan mengapa kecepatan pada sumbu yang diselesaikan perlu dinolkan.
        Sebutkan gejala yang muncul apabila langkah itu dilupakan.

  \item Sebuah permainan memuat benda berbentuk bintang yang berputar.
        Tentukan metode pemeriksaan tabrakan yang tepat untuk benda itu,
        beserta alasan dan pertimbangan biaya perhitungannya.
\end{enumerate}

\begin{refleksibox}
\begin{enumerate}
  \item Pada langkah ketiga praktikum, gejala apa yang pertama kali Anda
        sadari? Apakah Anda menemukan penyebabnya sendiri, atau perlu
        dibandingkan dengan program semula?
  \item Program bab ini memeriksa kebenaran tabrakannya
        sendiri melalui penghitung tumpang tindih. Adakah bagian lain pada
        program Anda yang menurut Anda juga layak diperiksa sendiri?
  \item Blok~I telah selesai. Bagian mana yang menurut Anda paling menantang,
        dan bagian mana yang paling berguna bagi pekerjaan Anda nanti?
\end{enumerate}
\end{refleksibox}
